% Research note prepared during checkpoint05 freeze; not yet included in paper.
\begin{saresult}{D.026}{Proved}{All remaining quadruple automorphism groups}
For ordered apposition quadruples in a simple adjoint complex group,
the following list, together with D.025, determines every possible
automorphism group and the number of inner-conjugacy classes of each type.
\begin{itemize}
\item In $B_r,C_r,G_2,F_4,E_7,E_8$, every automorphism group is trivial.
\item In $D_r$, $r\ge5$, put
$N_r=|W(D_r)|^2/(8(r-1)\varphi(2r-2))$.
Exactly $2N_r/r^2$ classes have automorphism group $C_2$;
the other $N_r-2N_r/r^2$ have trivial automorphism group.
\item In $D_4$, among the $768$ classes, $504$ have trivial
automorphism group, $252$ have group $C_2$, and $12$ have group $S_3$.
None has group $C_3$.  The quotient by all automorphisms has $180$ classes.
\item In $E_6$, among the $18662400$ classes, $27648$ have group
$C_2$ and $18634752$ have trivial automorphism group.
The quotient by all automorphisms has $9345024$ classes.
\end{itemize}
The finite certificates also determine the stabilizer of an individual
class by the exact twisted-Weyl criterion in the proof.
\end{saresult}
\begin{proof}
\medskip\noindent\textbf{Equivariant identification of the two Weyl torsors.}
Let $\Gamma$ be the pinned diagram-automorphism group.
In types $D_r,E_6$, the cyclic element
$x=\sum_i x_{\alpha_i}+x_{-\theta}$ and the principal element
$t=\rho^\vee(\zeta)$ can be chosen fixed by $\Gamma$; the sign
calculation D.005 verifies this on the highest-root vector.
Thus $T$ and $T'=C_G(x)$ are $\Gamma$-stable and $\Gamma$ fixes
the regular semisimple vector $x\in\operatorname{Lie}T'$.
The fixed subspace of the real cocharacter space of $T'$ contains
a regular vector: otherwise it would be contained in some root
hyperplane, and its complexification could not contain $x$.
More explicitly, finitely many proper real root hyperplanes
cannot cover that real vector space.  Choose its chamber;
the associated Borel $B'_0$ is $\Gamma$-stable.
Together with the pinned Borel $B_0\supset T$, this identifies
the two Borel sets with Weyl groups on which $\Gamma$ acts by
diagram automorphisms.  The action on either diagram has the
same outer class as its action on $G$, and is faithful.

The Coxeter element $c$ induced by $t$ on $T'$ commutes with
$\Gamma$.  We may identify it with the following standard
$\Gamma$-invariant Coxeter element.  In $D_r$ with the end-node
involution, use a negative $(r-1)$-cycle and a negative singleton
in signed coordinates.  Every invariant element in the $D_r$
Coxeter conjugacy class is conjugate to this one by the fixed
Weyl subgroup: commuting with the sign change on the final
coordinate forces the unique singleton of the signed cycle type
to occupy that coordinate.  The remaining negative cycle is
unique up to signed permutation of the other coordinates.
For full triality in $D_4$, take
$c=s_1s_3s_4s_2$ in Bourbaki numbering.  Complete enumeration
finds precisely two triality-fixed elements of its Weyl conjugacy
class, and they are a single orbit under $W^{S_3}$.
For $E_6$, take the bipartite product
$c=(s_1s_4s_6)(s_2s_3s_5)$.
Complete enumeration finds precisely $96$ diagram-fixed elements
of its Weyl conjugacy class, and they are a single orbit under
$W^{C_2}$; these cardinalities are certified, not assumed from
a folding picture.  This establishes the equivariant comparison
needed to use the following finite model for the geometric action.

Write $A=N_W(\langle c\rangle)$ and
$B_t=\{b\in W:b(t)=t^a,\ a\in U_h\}$.
By D.003 the common normalizer is
$\Phi=B_t\times_{U_h}A$, acting freely on $W\times W$ by
left multiplication in the two factors.  The outer group acts
by $(w,w')\mapsto(\gamma(w),\gamma(w'))$.
The stabilizer of the class of $(w,w')$ consists exactly of
those $\gamma\in\Gamma$ for which
\[
 \bigl(\gamma(w)w^{-1},\gamma(w')w'^{-1}\bigr)\in\Phi.
                                                               \tag{D26.1}
\]
As the inner stabilizer is trivial, this subgroup is precisely
the automorphism group of the corresponding quadruple.
This also supplies the individual-class test.

For an involution $\sigma$, define
$\Delta_\sigma=\{\sigma(w)w^{-1}:w\in W\}$.
Every fibre of $w\mapsto\sigma(w)w^{-1}$ has cardinality
$|W^\sigma|$: equality for $w_1,w_2$ is equivalent to
$w_2^{-1}w_1\in W^\sigma$.
Counting (D26.1) before dividing by the free $\Phi$-action gives
\[
 \#\operatorname{Fix}(\sigma;\Phi\backslash(W\times W))
 =\frac{|W^\sigma|^2}{|\Phi|}
   \#\bigl(\Phi\cap(\Delta_\sigma\times\Delta_\sigma)\bigr).
                                                               \tag{D26.2}
\]
This counts every point of the finite quotient, even when that
quotient is much larger than $W$.

\medskip\noindent\textbf{The uniform $D_r$ calculation.}
Put $m=r-1\ge3$, so $h=2m$ and $z=4$.
Use signed-coordinate matrices.  The end-node involution is
conjugation by the sign change in the last coordinate.
Its coboundary set consists exactly of the identity and the
$r-1$ diagonal matrices changing signs at coordinate $j$ and
the last coordinate.  Also
$|W^\sigma|=2^{r-1}(r-1)!$.
Let $\rho^\vee=(m,m-1,\ldots,1,0)$.
A nontrivial such diagonal matrix belongs to $B_t$ only if
it changes the coordinate with value $m$; its exponent is then one.
Indeed the zero last coordinate excludes the half-integral
coset of $hP^\vee$, and an unchanged coordinate of value one,
or unchanged coordinates of values two and three, force
$a=1\pmod {2m}$.  The changed value must then be a multiple of $m$.
Thus $B_t\cap\Delta_\sigma$ has two elements, both of exponent one.

On the other hand $A\cap\Delta_\sigma=\{1\}$.
Conjugating the negative $m$-cycle by one coordinate sign change
alters precisely two of its edge signs.  Its underlying permutation
is unchanged, so equality with $c^a$ requires $a\equiv1\pmod m$.
The only alternative to $a=1$ is $a=m+1$, available only for even
$m$; it changes all $m$ edge signs.  Since $m\ge3$, neither
alternative equals the change of exactly two edge signs.
Consequently the intersection in (D26.2) has two elements.
Formula (D26.2) gives
$2|W^\sigma|^2/|\Phi|=2N_r/r^2$.
For $r\ge5$ the outer group is $C_2$, which proves the asserted
full classification.  The exact finite checks for $4\le r\le20$
are in \path{code/type_d_outer_stabilizers.py}; the uniform argument,
not extrapolation, proves the statement for every rank.

\medskip\noindent\textbf{Triality in $D_4$.}
The program enumerates all $192$ Weyl matrices, all $12$ elements
of $A$, all $8$ elements of $B_t$, and all $48$ elements of $\Phi$.
It then enumerates all $768$ free orbits in $W\times W$ and
checks all six diagram automorphisms on each orbit.
The stabilizer orders occur with multiplicities
\[
\begin{array}{c|rrrr}
\text{order}&1&2&3&6\\\hline
\text{number of classes}&504&252&0&12.
\end{array}
\]
As an independent count, for each $\gamma\in S_3$ the program
sums over $(b,a)\in\Phi$ the product of the solution counts
$\gamma(w)=bw$ and $\gamma(w')=aw'$.
After division by $48$, the fixed-class counts are
$768$ for the identity, $96$ for each transposition, and $12$
for each three-cycle.  These agree with the complete class list.
They give $180$ full-automorphism orbits by averaging.
The matrices, class representatives, and actual stabilizer subgroups
are all in \path{logs/d4_outer_stabilizers.json}; the two methods
and the Coxeter-class comparison run in the accompanying program.

\medskip\noindent\textbf{The $E_6$ calculation.}
The program enumerates all $51840$ Weyl matrices.
It finds $|W^\sigma|=1152$ and exactly $45$ coboundaries,
with every fibre counted independently to have size $1152$.
The sets $A,B_t,\Phi$ have cardinalities $48,12,144$.
Exactly three pairs in $\Phi$ have both entries in
$\Delta_\sigma$; their second entries are all the identity.
Thus (D26.2) gives
\[
 \frac{3\cdot1152^2}{144}=27648.
\]
Subtracting from $51840^2/144=18662400$ gives the stated trivial
stabilizer count.  Every tested matrix, every coboundary with a
witness, the three contributing pairs, and the fixed-Coxeter
conjugacy check are saved in \path{logs/e6_outer_stabilizers.json}.
The complete run, including the geometric-comparison certificate,
takes $4.801$ seconds.  The other simple types have trivial outer
group, so D.004 already makes their quadruple automorphism groups
trivial.  This completes all simple complex types.
\end{proof}
